% !TeX root = ./main.tex

\xynusetup{
  title                = {\texorpdfstring{\CJKunderline[thickness=1.5bp, depth=0.12em, sep=0.04em]{信阳师范大学学位论文模板示例文档 v0.1}}{信阳师范大学学位论文模板示例文档 v0.1}},
  title*               = {An example of thesis template for Xinyang Normal University v0.1},
  author               = {本人姓名},
  author*              = {Li Feiyu},
  speciality           = {所学专业名称},
  speciality*          = {Mathematics and Applied Mathematics},
  supervisor           = {XXX（副）教授},% 导师
  supervisor*          = {Prof. XXX},
  practice-supervisor  = {XXX~教授},  % 校外导师
  practice-supervisor* = {XXX},
  % degree-category      = {教育硕士},  % 申请学位类别（默认按 degree-type 自动填写）
  % date                 = {2026-01-01},  % 论文提交日期，默认为今日
   department           = {学院名称},  % 所属院（系、所）
   department* = {School of Foreign Language},
  % classification-number = {TU473},     % 分类号（如需）
  % school-code           = {},    % 学校代码（需用户自行设置）
  % student-id           = {9527},  % 学号（如需）
  % secret-level         = {秘密},     % 绝密|机密|秘密|
  % secret-level*        = {Secret},  % Top secret | Highly secret | Secret
  % 数学字体
  % math-style           = GB,  % 可选：GB, TeX, ISO
  math-font            = xits,  % 可选：stix, xits, libertinus
}


% 加载宏包

% 定理类环境宏包
\usepackage{amsthm}

% 插图
\usepackage{graphicx}
\usepackage{float}
% CJK 可断行下划线（适用于 XeLaTeX）
\usepackage{xeCJKfntef}

% 英文图题、表题
\usepackage{bicaption}

% 分图
\usepackage{subcaption}

% 三线表
\usepackage{booktabs}

% 表注
\usepackage{threeparttable}

% 跨页表格
\usepackage{longtable}

% 算法
\usepackage[ruled,linesnumbered]{algorithm2e}

% SI 量和单位
\usepackage{siunitx}

% 参考文献使用 BibTeX + natbib 宏包
% 顺序编码制
\usepackage[sort]{natbib}
\bibliographystyle{xynuthesis-numerical}

% 著者-出版年制
% \usepackage{natbib}
% \bibliographystyle{xynuthesis-authoryear}

% 参考文献使用 BibLaTeX 宏包
% \usepackage[style=xynuthesis-numeric]{biblatex}
% \usepackage[bibstyle=xynuthesis-numeric,citestyle=xynuthesis-inline]{biblatex}
% \usepackage[style=xynuthesis-authoryear]{biblatex}
% 声明 BibLaTeX 的数据库
% \addbibresource{bib/xynu.bib}

% 配置图片的默认目录
\graphicspath{{figures/}}

% 数学命令
\makeatletter
\newcommand\dif{%  % 微分符号
  \mathop{}\!%
  \ifxynu@math@style@TeX
    d%
  \else
    \mathrm{d}%
  \fi
}
\makeatother
\newcommand\eu{{\symup{e}}}
\newcommand\iu{{\symup{i}}}

% 用于写文档的命令
\DeclareRobustCommand\cs[1]{\texttt{\char`\\#1}}
\DeclareRobustCommand\env[1]{\texttt{#1}}
\DeclareRobustCommand\pkg[1]{\textsf{#1}}
\DeclareRobustCommand\file[1]{\nolinkurl{#1}}

% 引入生成英文乱数假文的宏包
\usepackage{lipsum} 
% 引入生成中文乱数假文的宏包
\usepackage{zhlipsum}

% hyperref 宏包在最后调用
\usepackage{hyperref}
